\chapter{Metodolog\'ia}

\label{chap:implementacion}

En este cap\'itulo se describe la implementaci\'on computacional desarrollada para replicar y adaptar la metodolog\'ia propuesta por Kurt~\cite{Kurt2024}, en la cual se eval\'ua el desempe~no de motores de aeronaves durante la fase de despegue mediante t\'ecnicas de aprendizaje autom\'atico. A diferencia del trabajo original, que emplea el entorno MATLAB, la presente implementaci\'on se desarroll\'o en Python, utilizando la librer\'ia scikit-learn para el entrenamiento y evaluaci\'on de los modelos de Regresi\'on de Procesos Gaussianos (GPR), M\'aquinas de Vectores de Soporte (SVM) y Perceptr\'on Multicapa (MLP).

El conjunto de datos corresponde al \textit{ICAO Engine Emission Databank} (EEDB). Se emplean las siete variables de entrada identificadas en el marco te\'orico: tipo de motor, relaci\'on de derivaci\'on, relaci\'on de presi\'on, empuje nominal, presi\'on bar\'ometrica ambiente, temperatura ambiente y humedad ambiente. La variable objetivo es el par\'ametro \textbf{Fuel Flow T/O (kg/sec)}.

\section{Importaci\'on de librer\'ias y configuraci\'on}

El primer paso consiste en importar las librer\'ias necesarias para el manejo de datos (\texttt{numpy}, \texttt{pandas}), la visualizaci\'on (\texttt{matplotlib}) y los modelos de aprendizaje autom\'atico (\texttt{scikit-learn}). Se fija una semilla aleatoria (\texttt{RANDOM\_STATE = 42}) para garantizar la reproducibilidad de los resultados.

\begin{lstlisting}[language=Python, caption={Importaciones y configuracion inicial}, label={lst:imports}]
import numpy as np
import pandas as pd
import matplotlib.pyplot as plt
import matplotlib.gridspec as gridspec
import warnings
warnings.filterwarnings('ignore')

from sklearn.svm import SVR
from sklearn.neural_network import MLPRegressor
from sklearn.gaussian_process import GaussianProcessRegressor
from sklearn.gaussian_process.kernels import RBF, RationalQuadratic, Matern
from sklearn.preprocessing import StandardScaler
from sklearn.model_selection import train_test_split
from sklearn.metrics import mean_squared_error, mean_absolute_error

RANDOM_STATE = 42
np.random.seed(RANDOM_STATE)
\end{lstlisting}

\section{Carga y preparaci\'on del dataset}

A partir del archivo EEDB se construyen las variables de entrada mediante un proceso de ingenier\'ia de caracter\'isticas consistente con la metodolog\'ia del art\'iculo. La variable \texttt{Engine Type} se codifica de forma binaria (turbofan simple = 1, turbofan mezclado/\textit{mixed turbofan} = 2), mientras que las variables ambientales (presi\'on bar\'ometrica, temperatura y humedad) se calculan como el promedio entre sus valores m\'inimo y m\'aximo reportados en la base de datos, dado que el art\'iculo original reporta un \'unico valor representativo para cada una.

\begin{lstlisting}[language=Python, caption={Carga y preparacion del dataset EEDB}, label={lst:dataset}]
EXCEL_PATH = 'edb-emissions-databank_v31__web__0.xlsx'
df_raw = pd.read_excel(EXCEL_PATH, sheet_name='Gaseous Emissions and Smoke')

# Ingenier'ia de features segun el articulo
df_raw['Engine_Type'] = df_raw['Eng Type'].apply(
    lambda x: 1 if str(x).strip() == 'TF' else 2)

df_raw['Ambient_Baro'] = (df_raw['Ambient Baro Min (kPa)']
                           + df_raw['Ambient Baro Max (kPa)']) / 2
df_raw['Ambient_Temp'] = (df_raw['Ambient Temp Min (K)']
                           + df_raw['Ambient Temp Max (K)']) / 2
df_raw['Ambient_Humidity'] = (df_raw['Humidity Min (kg/kg)']
                               + df_raw['Humidity Max (kg/kg)']) / 2

FEATURES = ['Engine_Type', 'B/P Ratio', 'Pressure Ratio', 'Rated Thrust (kN)',
            'Ambient_Baro', 'Ambient_Temp', 'Ambient_Humidity']
TARGET   = 'Fuel Flow T/O (kg/sec)'

df = df_raw[FEATURES + [TARGET]].dropna().reset_index(drop=True)
\end{lstlisting}

Tras eliminar los registros con valores faltantes, el conjunto depurado queda conformado por motores certificados descritos por las siete variables predictoras y el valor real de \textit{Fuel Flow T/O}.

\noindent\textbf{N\'umero de motores usados.} El art\'iculo de referencia usa dos ediciones del EEDB: la edici\'on 09/2019 con 565 registros, empleados para entrenar, validar y probar los tres modelos, y la edici\'on 07/2021 con 773 registros, empleados s\'olo al final para probar el mejor modelo con datos nuevos. En total se usan 1 338 registros (565 + 773). La r\'eplica de este proyecto, detallada en REPLICACION.md y construida en el c\'odigo de \texttt{src/} con NumPy puro, obtiene 563 registros para 2019 tras excluir filas incompletas y 773 para 2021, es decir 1 336 en total, lo cual coincide con el art\'iculo.

El texto original de esta memoria indicaba 813 motores. Ese n\'umero no coincide con los datos que se usan de verdad. El c\'odigo de la tesis en scikit-learn lee un archivo distinto del EEDB (\texttt{edb-emissions-databank\_v31\_\_web\_\_0.xlsx}) y, tras eliminar registros con valores faltantes, queda con 813 motores. Esa cifra corresponde a una edici\'on distinta de la base de datos y no coincide ni con el art\'iculo (565 y 773) ni con la implementaci\'on de \texttt{src/}, que usa las dos ediciones citadas y NumPy puro. Por lo tanto, 813 es un dato del archivo v3.1 scikit-learn y no del conjunto de estudio.

\section{An\'alisis de regresi\'on m\u00faltipla}

Antes de entrenar los modelos de aprendizaje autom\'atico, se realiza un an\'alisis de regresi\'on lineal m\u00faltipla con el fin de verificar la relaci\'on estad\'istica entre las variables de entrada y el par\'ametro objetivo, replicando el an\'alisis de la Secci\'on 2.2 del art\'iculo de referencia.

\begin{lstlisting}[language=Python, caption={Regresion lineal m\u00faltipla de referencia}, label={lst:regresion}]
from sklearn.linear_model import LinearRegression

X_reg = df[FEATURES].values
y_reg = df[TARGET].values

lr = LinearRegression()
lr.fit(X_reg, y_reg)

y_pred_lr = lr.predict(X_reg)
r2 = 1 - np.sum((y_reg - y_pred_lr)**2) / np.sum((y_reg - y_reg.mean())**2)
r  = np.sqrt(r2)
\end{lstlisting}

El coeficiente de correlaci\'on obtenido ($R$) confirma que existe una fuerte relaci\'on lineal entre las variables seleccionadas y el \textit{Fuel Flow T/O}, justificando su uso como entradas de los modelos no lineales posteriores.

\section{Justificaci\'on de las siete variables de entrada}

Tres de las siete variables presentan resultados at\'ipicos en el an\'alisis:

\begin{itemize}
    \item \textit{Press ratio}: Sig. $= 0.963$ ($t = -0.046$). Por s\'i sola, la variable no es significativa. (Tabla 9)
    \item \textit{Ambient baro}: m\'in. $= 0.00$ y m\'ax. $= 1\,019\,394.5$. La presi\'on real a nivel del suelo est\'a entre 95 y 104 kPa. (Tabla 4)
    \item \textit{Ambient temp}: m\'in. $= 144$ K ($-129$~$^{\circ}$C). El EEDB registra 288 K para ese motor (4AL002); el valor est\'a dividido entre 2. (Tablas 4 y 6)
\end{itemize}

Adem\'as, la Tabla 9 publica $\beta = -0.49$ para el tipo de motor, con error t\'ipico 0.017 y $t = -2.962$. Esos dos valores dan $\beta = -0.049$; el $-0.49$ es un error de un factor 10.

Aunque estas variables tienen resultados at\'ipicos, se conservan las siete por tres razones t\'ecnicas:

\begin{enumerate}
    \item La regresi\'on lineal s\'olo mide efectos lineales. \textit{Press ratio} no tiene efecto lineal propio, pero puede aportar informaci\'on a un modelo no lineal (GPR, SVM, MLP).
    \item La prueba ANOVA eval\'ua el conjunto de variables, no cada variable sola. El conjunto es significativo (Sig. $< 0.001$).
    \item Las siete variables son par\'ametros f\'isicos del motor y del ensayo de certificaci\'on. Todas tienen sentido f\'isico para el consumo de combustible.
\end{enumerate}

Con los datos corregidos del EEDB, la regresi\'on da el mismo resultado: $R^2 = 0.980$ y $t \approx 0$ para \textit{press ratio}. Los valores at\'ipicos no cambian la conclusi\'on de la regresi\'on.

\section{Divisi\'on del conjunto de datos y normalizaci\'on}

Siguiendo el mismo esquema empleado en el art\'iculo, el conjunto de datos se divide en tres subconjuntos: 70 \% para entrenamiento, 15 \% para validaci\'on y 15 \% para prueba. Adicionalmente, se aplica una normalizaci\'on est\'andar (\texttt{StandardScaler}) ajustada \'unicamente sobre el conjunto de entrenamiento, para evitar fuga de informaci\'on (\textit{data leakage}) hacia los conjuntos de validaci\'on y prueba.

\begin{lstlisting}[language=Python, caption={Division del dataset y normalizacion}, label={lst:split}]
X = df[FEATURES].values
y = df[TARGET].values

X_trainval, X_test, y_trainval, y_test = train_test_split(
    X, y, test_size=0.15, random_state=RANDOM_STATE)
X_train, X_val, y_train, y_val = train_test_split(
    X_trainval, y_trainval, test_size=0.15/0.85, random_state=RANDOM_STATE)

scaler = StandardScaler()
X_train_s = scaler.fit_transform(X_train)
X_val_s   = scaler.transform(X_val)
X_test_s  = scaler.transform(X_test)
\end{lstlisting}

\section{M\'etricas de evaluaci\'on}

Con el fin de comparar el desempe~no de los distintos modelos bajo un criterio uniforme, se implementa una funci\'on auxiliar que calcula las m\'etricas de error definidas en el marco te\'orico: el coeficiente de correlaci\'on ($R$), el error cuadr\'atico medio (MSE), el error absoluto medio (MAE) y el error porcentual absoluto medio (MAPE).

\begin{lstlisting}[language=Python, caption={Funcion de calculo de metricas de error}, label={lst:metricas}]
def compute_metrics(y_true, y_pred):
    mse  = mean_squared_error(y_true, y_pred)
    mae  = mean_absolute_error(y_true, y_pred)
    mape = np.mean(np.abs((y_true - y_pred) / y_true)) * 100
    r    = np.corrcoef(y_true, y_pred)[0, 1]
    return dict(R=r, MSE=mse, MAE=mae, MAPE=mape)
\end{lstlisting}

Estas m\'etricas se calculan de forma independiente para los conjuntos de entrenamiento, validaci\'on y prueba, lo que permite identificar posibles casos de sobreajuste (\textit{overfitting}) al comparar el error entre subconjuntos.

\section{Regresi\'on de Procesos Gaussianos (GPR)}

El primer modelo evaluado corresponde a la Regresi\'on por Proceso Gaussiano. Se prueban tres funciones de kernel equivalentes a las reportadas en el art\'iculo: el kernel exponencial (aproximado mediante un kernel de Mat\'ern con $\nu = 0.5$), el kernel cuadr\'atico racional y el kernel exponencial cuadrado (RBF).

\begin{lstlisting}[language=Python, caption={Entrenamiento de los modelos GPR}, label={lst:gpr}]
gpr_kernels = {
    'Exponential GPR'        : Matern(nu=0.5),
    'Rational Quadratic GPR' : RationalQuadratic(),
    'Squared Exponential GPR': RBF(),
}

gpr_models  = {}
gpr_metrics = {}

for name, kernel in gpr_kernels.items():
    gpr = GaussianProcessRegressor(kernel=kernel, n_restarts_optimizer=5,
                                    alpha=1e-6, random_state=RANDOM_STATE)
    gpr.fit(X_train_s, y_train)
    m_tr = compute_metrics(y_train, gpr.predict(X_train_s))
    m_va = compute_metrics(y_val,   gpr.predict(X_val_s))
    m_te = compute_metrics(y_test,  gpr.predict(X_test_s))
    gpr_models[name]  = gpr
    gpr_metrics[name] = {'train': m_tr, 'val': m_va, 'test': m_te}

best_gpr = min(gpr_metrics, key=lambda k: gpr_metrics[k]['test']['MSE'])
\end{lstlisting}

El mejor modelo GPR se selecciona autom\'aticamente como aquel que minimiza el MSE sobre el conjunto de prueba.

\section{M\'aquina de Vectores de Soporte (SVM)}

El segundo modelo corresponde a una M\'aquina de Vector Soporte para regresi\'on (SVR). Se eval\'uan tres variantes de kernel: lineal, cuadr\'atico (polinomio de grado 2) y c\'ubico (polinomio de grado 3), donde los hiperpar\'ametros del SVM cuadr\'atico (\texttt{C = 1.2076}, \texttt{epsilon = 0.1208}) se fijan de acuerdo con los valores estructurales reportados en la Tabla~13 del art\'iculo original.

\begin{lstlisting}[language=Python, caption={Entrenamiento de los modelos SVM}, label={lst:svm}]
svm_configs = {
    'Linear SVM'   : SVR(kernel='linear', C=1.0, epsilon=0.1),
    'Quadratic SVM': SVR(kernel='poly', degree=2, C=1.2076, epsilon=0.1208),
    'Cubic SVM'    : SVR(kernel='poly', degree=3, C=1.0, epsilon=0.1),
}

svm_models  = {}
svm_metrics = {}

for name, svm in svm_configs.items():
    svm.fit(X_train_s, y_train)
    m_tr = compute_metrics(y_train, svm.predict(X_train_s))
    m_va = compute_metrics(y_val,   svm.predict(X_val_s))
    m_te = compute_metrics(y_test,  svm.predict(X_test_s))
    svm_models[name]  = svm
    svm_metrics[name] = {'train': m_tr, 'val': m_va, 'test': m_te}

best_svm = min(svm_metrics, key=lambda k: svm_metrics[k]['test']['MSE'])
\end{lstlisting}

\section{Perceptr\'on Multicapa (MLP)}

El tercer modelo replica la arquitectura de red neuronal artificial descrita en la Figura~5 y la Tabla~15 del art\'iculo (arquitectura 10-8-4-3 con funciones de activaci\'on \textit{tansig}/\textit{logsig}). Dado que \texttt{scikit-learn} no implementa directamente el algoritmo de entrenamiento Levenberg--Marquardt (\texttt{trainlm}) utilizado en MATLAB, este se aproxima mediante el solucionador \texttt{lbfgs}, un m\'etodo cuasi-Newton de comportamiento equivalente para redes de tama~no moderado. De forma an\'aloga, el algoritmo de retropropagaci\'on resiliente (\texttt{trainrp}) se aproxima mediante el solucionador \texttt{sgd} con tasa de aprendizaje adaptativa.

\begin{lstlisting}[language=Python, caption={Entrenamiento de los modelos MLP}, label={lst:mlp}]
mlp_configs = [
    ('Model 1 trainlm (10-8-4-3)', MLPRegressor(
        hidden_layer_sizes=(10,8,4,3), activation='tanh', solver='lbfgs',
        max_iter=3000, random_state=RANDOM_STATE)),
    ('Model 2 trainlm (15-10-5)',   MLPRegressor(
        hidden_layer_sizes=(15,10,5),  activation='tanh', solver='lbfgs',
        max_iter=3000, random_state=RANDOM_STATE)),
    ('Model 3 trainlm (20-10)',     MLPRegressor(
        hidden_layer_sizes=(20,10),    activation='tanh', solver='lbfgs',
        max_iter=3000, random_state=RANDOM_STATE)),
    ('Model 4 trainrp (10-8-4-3)',  MLPRegressor(
        hidden_layer_sizes=(10,8,4,3), activation='tanh', solver='sgd',
        learning_rate='adaptive', max_iter=3000, random_state=RANDOM_STATE)),
    ('Model 5 trainrp (15-10-5)',   MLPRegressor(
        hidden_layer_sizes=(15,10,5),  activation='tanh', solver='sgd',
        learning_rate='adaptive', max_iter=3000, random_state=RANDOM_STATE)),
]

mlp_models  = {}
mlp_metrics = {}

for name, mlp in mlp_configs:
    mlp.fit(X_train_s, y_train)
    m_tr = compute_metrics(y_train, mlp.predict(X_train_s))
    m_va = compute_metrics(y_val,   mlp.predict(X_val_s))
    m_te = compute_metrics(y_test,  mlp.predict(X_test_s))
    mlp_models[name]  = mlp
    mlp_metrics[name] = {'train': m_tr, 'val': m_va, 'test': m_te}

best_mlp = min(mlp_metrics, key=lambda k: mlp_metrics[k]['test']['MSE'])
\end{lstlisting}

Se eval\'uan cinco configuraciones distintas, variando tanto la arquitectura de capas ocultas como el algoritmo de optimizaci\'on, con el fin de identificar la combinaci\'on que minimiza el error en el conjunto de prueba, de forma an\'aloga a la comparaci\'on presentada en la Tabla~14 del art\'iculo original.

\section{Comparaci\'on de modelos}

Una vez seleccionado el mejor modelo de cada familia (GPR, SVM y MLP) seg\'un el criterio de menor MSE en el conjunto de prueba, se construye una tabla comparativa con las m\'etricas MSE, MAE y MAPE para los tres conjuntos de datos (entrenamiento, validaci\'on y prueba), replicando la estructura de la Tabla~16 del art\'iculo original.

\begin{lstlisting}[language=Python, caption={Comparacion final entre modelos}, label={lst:comparacion}]
rows = []
for split_name, key in [('Train', 'train'), ('Val', 'val'), ('Test', 'test')]:
    for model_label, met_dict in [
        ('Exponential GPR', gpr_metrics[best_gpr]),
        ('MLP (10-8-4-3)',   mlp_metrics[best_mlp]),
        ('Linear SVM',       svm_metrics[best_svm]),
    ]:
        m = met_dict[key]
        rows.append({'Split': split_name, 'Modelo': model_label,
                     'MSE': m['MSE'], 'MAE': m['MAE'], 'MAPE': m['MAPE']})

results_df = pd.DataFrame(rows)
\end{lstlisting}

Adicionalmente, se genera una visualizaci\'on comparativa del MAPE de los tres modelos en cada conjunto de datos, incorporando una l\'inea de referencia en 10\%, umbral que la literatura (clasificaci\'on de Lewis, 1982) considera como el l\'imite superior para catalogar un modelo como de alta precisi\'on.

\section{Construcci\'on de intervalos de confianza}

Finalmente, se implementa el mecanismo de diagn\'ostico propuesto en la Secci\'on 3.5 del art\'iculo: a partir de las predicciones del mejor modelo (el MLP), se calculan los residuos sobre la totalidad del conjunto de datos y se construye un intervalo de confianza al 99\% ($z_{0.005} = 2.58$), definido como:

\begin{equation}
\text{Fuel Flow}_{T/O} \pm z_{\alpha/2} \cdot \sigma
\label{eq:intervalo_confianza}
\end{equation}

donde $\sigma$ es la desviaci\'on est\'andar de los residuos del modelo. Todo motor cuyo valor real de \textit{Fuel Flow T/O} se encuentre fuera de este intervalo se marca como una posible anomal\'ia de desempe~no.

\begin{lstlisting}[language=Python, caption={Calculo de intervalos de confianza al 99\% y deteccion de anomal'ias}, label={lst:confianza}]
best_model = mlp_models[best_mlp]

X_all_s  = scaler.transform(X)
ff_pred  = best_model.predict(X_all_s)

z_99      = 2.58
residuals = y - best_model.predict(X_all_s)
sigma     = np.std(residuals)

upper = ff_pred + z_99 * sigma
lower = ff_pred - z_99 * sigma

outliers = (y < lower) | (y > upper)
\end{lstlisting}

Este procedimiento constituye la aportaci\'on central del modelo: permite evaluar autom\'aticamente, sin c\'alculos manuales ni conocimiento experto, si el motor de una aeronave opera dentro de los par\'ametros esperados durante la fase de despegue, o si presenta indicios de deterioro de desempe~no.

\section{Resumen de la implementaci\'on}

La Tabla~\ref{tab:resumen_implementacion} sintetiza la correspondencia entre las secciones del art\'iculo original y los bloques implementados en el presente trabajo.

\begin{table}[H]
\centering
\caption{Correspondencia entre la metodolog\'ia del art\'iculo y la implementaci\'a desarrollada}
\label{tab:resumen_implementacion}
\begin{tabular}{@{}p{3.8cm}p{3.8cm}p{5.6cm}@{}}
\toprule
\textbf{Secci´on del art´iculo} & \textbf{Herramienta original (MATLAB)} & \textbf{Implementaci´on (Python)} \\
\midrule
Regresi´on m\u00faltipla (2.2)     & \texttt{regress}                    & \texttt{LinearRegression} \\
GPR (2.3.1 / 3.1)            & \textit{Regression Learner}          & \texttt{GaussianProcessRegressor} \\
SVM (2.3.2 / 3.2)            & \textit{Regression Learner}          & \texttt{SVR} \\
MLP (2.3.3 / 3.3)            & \textit{Neural Net Fitting}, \texttt{trainlm} & \texttt{MLPRegressor}, \texttt{solver='lbfgs'} \\
Intervalos de confianza (3.5) & C´alculo manual ($z\sigma$)          & \texttt{numpy} (Ec.~\ref{eq:intervalo_confianza}) \\
\bottomrule
\end{tabular}
\end{table}
